\section{Lean-Module, Übersetzung und Einsatz von Sprachmodellen}\label{app:lean}

\paragraph{Projektunabhängiges Paket.} \texttt{qam/lean/} ist ein eigenständiges Lake-Paket, das nur den Lean-Kern benutzt (kein Mathlib)
und in Sekunden übersetzt: \texttt{QAM/Provenance.lean} (Herleitungsbäume: keine Beförderung, Eingaben, Frames) und
\texttt{QAM/QamFile.lean} (der Prüfer für exportierte Akten, seine zeilen- und aktenweite Korrektheit, die Hash-Verweise) mit dem ausführbaren
Programm \texttt{qamcheck}.

\paragraph{Objektebene und Maschine.} Im Verzeichnis \texttt{formal/} des Repositorys formalisieren die Module der Bibliothek
\texttt{QSFormal} die Objektarten und die Maschine; diese Module benutzen ebenfalls nur den Lean-Kern, andere Module derselben Bibliothek
benutzen Mathlib:
\begin{itemize}[nosep]
\item \texttt{QSFormal/QAM.lean}: Disjunktheit der Markierungen (\texttt{qState\_iff\_tag}, \texttt{not\_state\_and\_op}, \texttt{tag\_lt\_eight}),
Umindizierung und Rückübersetzung (\texttt{back\_shift}, \texttt{shift\_back}), Transport und Ketten (\texttt{first\_transport\_principle},
\texttt{chain\_clEq}, \texttt{chain\_readout});
\item \texttt{QSFormal/Models.lean}: zwei triviale Modelle der Transportschemata (\texttt{totalModel}, \texttt{emptyModel}); die Schemata
sind widerspruchsfrei, legen aber nicht fest, was zulässig heißt;
\item \texttt{QSFormal/Route.lean}: Wege über mehrere Frames (\texttt{path\_admissible}, \texttt{path\_clEq}, \texttt{loop\_clEq});
\item \texttt{QSFormal/Store.lean}: Speicher, Anweisungen, Übergangssemantik, Erhaltung (\texttt{exec\_preserves}, \texttt{run\_preserves}),
Rücksetzen (\texttt{rollback\_restores}), Zertifikate (\texttt{no\_false\_certificate});
\item \texttt{QSFormal/StoreProgress.lean}: kein Steckenbleiben (\texttt{exec\_ne\_halt}), exakte Statuscharakterisierung, Atomarität
(\texttt{exec\_atomic}), Fortschritt für \textsc{snapshot}, \textsc{scrub} und \textsc{export}, Vollständigkeit des Protokolls.
\end{itemize}
Vereinfachungen gegenüber der vollen Spezifikation: Adressen sind natürliche Zahlen, Relationen tragen keine Zeigerbedingungen, und die
Anweisungen zum Kopieren und Zusammenführen von Speichern sind nicht formalisiert. Die kontinuierliche Integration übersetzt beide Teile
und weist \texttt{sorry}, \texttt{admit} und \texttt{native\_decide} zurück.

\paragraph{Dieser Artikel.} Der Artikel lädt \texttt{qam.sty}; die Beispiele in den Abschnitten~\ref{sec:latex} und~\ref{sec:fallstudie}
werden beim Setzen geprüft. Übersetzung im Verzeichnis \texttt{paper/qam}:
\begin{lstlisting}
TEXINPUTS=../../qam/tex: pdflatex main.tex   (zweimal)
\end{lstlisting}

\paragraph{Einsatz von Sprachmodellen.} Die Ausgangsvermutung der Fallstudie stammt vom Autor. Das Modell QAM, die Werkzeuge, die
Lean-Formalisierung und große Teile dieses Texts sind im Dialog mit großen Sprachmodellen entstanden (Claude, Anthropic; ChatGPT und Codex,
OpenAI). Manche Ideen kamen von den Modellen, andere vom Autor; der Autor hat die Fragen gewählt und entschieden, was bleibt. Ergebnisse der
Modelle wurden nur über die in Abschnitt~\ref{sec:fallstudie} beschriebenen Prüfungen übernommen. Die Verantwortung für den Inhalt trägt
der Autor.
